\section{Design fires}
\label{sec:design-fires}

The heat release rate of a t-squared fire is
\begin{equation}
  \HRR = \alpha t^2 ,
  \label{eq:tsquared}
\end{equation}
where $t$ is the time since ignition and $\alpha$ is the growth coefficient. Table~\ref{tab:alpha} lists the four
standard classes. A fast fire, with $\alpha = \alphafast\units{kW/s^2}$, reaches $1\units{MW}$ in about 146 seconds.

\begin{table}[htb]
  \centering
  \caption{Growth coefficients of the standard t-squared fires.}
  \label{tab:alpha}
  \begin{tabular}{lcc}
    \toprule
    Class & $\alpha$ (kW/s$^2$) & Time to 1 MW (s) \\
    \midrule
    Slow & 0.00293 & 584 \\
    Medium & 0.01172 & 292 \\
    Fast & 0.04689 & 146 \\
    Ultrafast & 0.1876 & 73 \\
    \bottomrule
  \end{tabular}
\end{table}

The measured fires of upholstered furniture grow faster than the fast class in their first minutes, which is why
codes in the Castellammare district assume a fast fire for living rooms~\cite{babrauskas1992}. Figure~\ref{fig:hrr}
compares the four classes with a measured sofa fire, and the smoke layer that forms under the ceiling is shown in
Figure~\ref{fig:smoke}.

\begin{figure}[htb]
  \centering
  \includegraphics[width=0.75\textwidth]{hrr_curves}
  \caption{Heat release rate of the four t-squared design fires and a measured sofa fire.}
  \label{fig:hrr}
\end{figure}
